% ====================================================================
% FILE: content/13_branching_topology.tex
% Онтологическое ветвление и глобальная топология
% Restored Full Version for Core 1.3.3
% ====================================================================

\section{Онтологическое ветвление и глобальная топология}
\label{sec:branching-topology}

В этом разделе реконструируется теория онтологических ветвей и глобальной
топологии континуумов. Он раскрывает аксиому~21 (онтологическое
ветвление) и связывает её с размерностной иерархией
\(K_0\)–\(K_{10}\), операторами рождения
\(\Psi_{x\to x+1}\) (раздел~the operator semantics chapter) и структурой
схлопывания\textendash возрождения (раздел~the lifecycle and rebirth chapter).

Неформально, \emph{ветвь} — это максимальный путь развития, вдоль которого
континуум последовательно подвергается размерностным рождениям, совместимым с
заданным набором осей и пространств вложения. Множество таких ветвей
несет естественную направленную топологию и даёт глобальный обзор всех
допустимых онтогенетических траекторий, совместимых с аксиомами OC.

\subsection{Определение онтологических ветвей}

Уровень континуума \(K_x\) задается:
\[
  K_x =
  \big(
    \Omega_x, A_x, P_x(t), J_x(t),
    \Theta_x, \partial\Omega_x, C_x, k_x(t)
  \big),
\]
вложенный в пространство \(M_x\) с множеством осей \(A(M_x)\). Размерностное
рождение из \(K_x\) в \(K_{x+1}\) опосредуется оператором рождения
\(\Psi_{x\to x+1}\) (раздел~the operator semantics chapter).

\paragraph{Определение 13.1 (Онтологическая ветвь).}
Онтологическая ветвь \(B\) — это конечная или счётная последовательность
континуумов
\[
  B = (K_{x_0}, K_{x_1}, K_{x_2}, \dots),
\]
вместе с пространствами вложения \((M_{x_0}, M_{x_1}, \dots)\), такая что:
\begin{enumerate}
  \item \textbf{Монотонность уровней.}\quad
        \(x_0 < x_1 < x_2 < \dots\) и каждый шаг
        \(K_{x_i} \to K_{x_{i+1}}\) представляет собой размерностный
        переход (не чисто параметрическую эволюцию);
  \item \textbf{Совместимость рождения.}\quad
        Для каждого \(i\) существует оператор рождения
        \(\Psi_{x_i\to x_{i+1}}\) такой, что
        \[
          (K_{x_{i+1}}, M_{x_{i+1}}) =
            \Psi_{x_i\to x_{i+1}}(K_{x_i}, M_{x_i}),
        \]
        где \(\Psi_{x_i\to x_{i+1}}\) удовлетворяет общим условиям рождения
        из раздела~the operator semantics chapter;
  \item \textbf{Монотонность осей.}\quad
        Множества осей растут монотонно:
        \[
          A_{x_i} \subset A_{x_{i+1}}
          \subseteq A(M_{x_{i+1}}),
        \]
        и каждая новая ось не представима в структурном спане предыдущих осей;
  \item \textbf{Монотонность вложения.}\quad
        Пространства вложения образуют цепочку включения
        \[
          M_{x_0} \subset M_{x_1} \subset M_{x_2} \subset \dots
        \]
        с \(A(M_{x_i}) \subseteq A(M_{x_{i+1}})\);
  \item \textbf{Максимальность.}\quad
        Ветвь максимальна относительно расширения в направлении увеличения
        уровня: не существует такого континуума \(K_{x^{\ast}}\) с
        \(x^{\ast} > \sup_i x_i\), для которого последовательность можно
        строго продолжить, сохраняя указанные выше условия.
\end{enumerate}

Если \(x_0 = 0\), ветвь имеет корень в структурном субстрате \(K_0\). Более
общо ветви могут определяться относительно любого исходного уровня
\(K_{x_0}\), выполняющего роль эффективного корня для ограниченной
предметной области (например, химии, когниции, социальных систем).

\paragraph{Замечание 13.2 (Идентичность ветвей).}
Две ветви \(B\) и \(B'\) считаются различными, если они различаются хотя бы
по одному из параметров:
\begin{itemize}
  \item последовательности уровней \(x_i\) (например, пропуск или вставка
        уровней);
  \item наборов осей \(A_{x_i}\) (разным возникшим степеням свободы);
  \item семейств порогов \(\Theta_{x_i}\) или пространств вложения \(M_{x_i}\).
\end{itemize}
Ветви таким образом кодируют \emph{истории развития} в пространстве
континуумов, а не просто статические наборы уровней.

\subsection{Топология ветвей}

Множество всех ветвей имеет естественную направленную структуру.

\paragraph{Определение 13.3 (Граф ветвей).}
Пусть \(\mathcal{K}\) — класс всех континуумов, допускаемых аксиомами OC, а
\(\mathcal{B}\) — множество всех ветвей. Тогда
\emph{граф ветвей} \(\mathcal{G}\) — это ориентированный граф, у которого:
\begin{itemize}
  \item вершины — континуумы \(K_x\in\mathcal{K}\);
  \item ориентированные рёбра \(K_x \to K_y\) соответствуют операторам
        рождения \(\Psi_{x\to y}\) с \(y > x\).
\end{itemize}
Ветвь \(B\) тогда является ориентированным путём в \(\mathcal{G}\),
максимальным относительно расширения.

Так как операторы рождения не уменьшают размерность (монотонность размерности),
\(\mathcal{G}\) является ацикличным ориентированным графом (DAG). Циклы
могут существовать в \emph{динамике состояний} данного континуума \(K_x\)
(через внутренние циклы \(C_x\)), но не в вертикальном графе рождения.

\paragraph{Глобальная топология.}
Глобальная топология \(\mathcal{G}\) ограничивается:
\begin{itemize}
  \item \textbf{Общими префиксами.}\quad
        Разные ветви имеют общие начальные сегменты, когда они совпадают на
        нижних уровнях (например, все физические, химические и
        биологические ветви разделяют \(K_0\)–\(K_2\));
  \item \textbf{Расхождением ветвей.}\quad
        Ветви расходятся, когда операторы рождения на некотором уровне
        \(K_x\) выбирают различные новые оси или пороговые структуры,
        приводящие к разным континуумам \(K_{x+1}^{(a)}\),
        \(K_{x+1}^{(b)}\) и т.п.;
  \item \textbf{Возможным слиянием.}\quad
        Через операторы взаимодействия \(E_{\mathrm{int}}\) разные ветви позже могут
        порождать континуумы, лежащие на общей расширенной траектории,
        что создает DAG-подобную, а не только древовидную топологию;
  \item \textbf{Гиперграфовой структурой.}\quad
        Когда операторы рождения принимают на вход несколько континуумов
        (например, события слияния), граф ветвей можно расширить до
        ориентированного гиперграфа с гиперрёбрами
        \((K_{x}^{(a)}, K_{x}^{(b)},\dots) \to K_{y}\).
\end{itemize}

Глобальная топология ветвей таким образом интерполирует между деревоподобными
структурами (чистая фиссия без слияний) и гиперграфоподобными структурами
(слияние ветвей и коллективное рождение). OC не фиксирует единую глобальную
форму, но ограничивает все допустимые формы с помощью монотонности
размерности, совместимости вложения и пороговых условий.

\subsection{Условия ветвления}

Ветвление не возникает в произвольных точках; оно структурно инициируется
специфическими пороговыми и вложенными условиями.

\paragraph{Аксиома 21 (онтологическое ветвление).}
Пусть \(K_x\) — континуум на уровне \(x\) со сферой вложения \(M_x\). Событие
ветвления в \(K_x\) допускается, если существуют два или более различных
оператора рождения
\[
  \Psi_{x\to x+1}^{(a)},\ \Psi_{x\to x+1}^{(b)},\dots
\]
таких, что:
\begin{enumerate}
  \item \textbf{Различные новые оси.}\quad
        Каждый оператор рождения выбирает структурно различную новую ось
        \(A_{\mathrm{new}}^{(a)}\), \(A_{\mathrm{new}}^{(b)}\), \dots,
        взаимно непредставимые в структурном спане, то есть
        \[
          A_{\mathrm{new}}^{(a)} \notin \mathrm{span}_{\mathrm{str}}
            \big(A_x \cup \{A_{\mathrm{new}}^{(b)}\}\big), \ \text{и т.д.};
        \]
  \item \textbf{Размерностное напряжение.}\quad
        Структурное напряжение превосходит размерностный порог
        по нескольким неэквивалентным направлениям различий:
        \[
          T(K_x,t) > \Theta_{\mathrm{dim}}^{(a)}(K_x)
          \quad\text{и}\quad
          T(K_x,t) > \Theta_{\mathrm{dim}}^{(b)}(K_x),
        \]
        где верхние индексы обозначают различные размерностные каналы;
  \item \textbf{Емкость вложения.}\quad
        Пространство вложения \(M_x\) поддерживает все кандидатные оси, то
        есть
        \[
          A_{\mathrm{new}}^{(a)}, A_{\mathrm{new}}^{(b)}, \dots
          \in A(M_x);
        \]
  \item \textbf{Непустые допустимые области.}\quad
        Для каждого кандидатного континуума
        \(K_{x+1}^{(a)} = \Psi_{x\to x+1}^{(a)}(K_x, M_x)\), и т.д.,
        допустимое множество непусто: \(\Omega_{x+1}^{(a)}\neq\emptyset\).
\end{enumerate}

Во многих конкретных системах данные условия реализуются через:
\begin{itemize}
  \item несогласованности в потенциалах \(P_x(t)\), которые нельзя
        устранить внутри одного продолжения;
  \item конкуренцию между комплексами циклов \(C_x^{(a)}\), \(C_x^{(b)}\),
        чья совместная стабилизация нарушила бы пороги устойчивости;
  \item геометрии границы, допускающей различные взаимно исключающие
        повторные замыкания (например, альтернативные архитектуры
        мембран или институциональные структуры).
\end{itemize}

\paragraph{Пороги ветвления.}
Для удобства можно определить \emph{порог ветвления}
\(\Theta_{\mathrm{branch}}(K_x)\) как эффективную композицию
размерностных и вложенных порогов, такую, что
\[
  T_{\mathrm{branch}}(K_x,t) > \Theta_{\mathrm{branch}}(K_x)
  \quad\Rightarrow\quad
  \text{ветвление возможно.}
\]
Core~1.3.3 рассматривает \(\Theta_{\mathrm{branch}}\) как производную величину,
а не как новый первичный класс порогов: она кодирует одновременную
выполнимость условий нескольких операторов рождения.

\subsection{Взаимодействие ветвей}

Ветви обычно не эволюционируют изолированно. Если континуумы из разных
ветвей сосуществуют в общем пространстве вложения, они могут взаимодействовать
через оператор взаимодействия \(E_{\mathrm{int}}\)
(раздел~the operator semantics chapter).

\paragraph{Определение 13.4 (Взаимодействие ветвей).}
Пусть \(B^{(a)}\) и \(B^{(b)}\) две ветви с континуумами
\(\{K_x^{(a)}\}\) и \(\{K_y^{(b)}\}\), соответственно, вложенные в
общее или пересекающееся пространство \(M\). \emph{Взаимодействие ветвей} —
семейство операторов взаимодействия
\[
\begin{aligned}
  E_{\mathrm{int}}^{x,y}:\quad
    &(K_x^{(a)}, K_y^{(b)}, M)\\
    &\longrightarrow (K_x^{(a)\prime}, K_y^{(b)\prime}, M'),
\end{aligned}
\]
определенное для релевантных пар уровней \((x,y)\).

Далее перечислены типичные режимы:

\begin{itemize}
  \item \textbf{Сосуществование.}\quad
        Потоки слабо связаны и позволяют обеим ветвям удерживать или увеличивать
        свою континуумичность:
        \[
          k^{(a)}(t+dt) \ge k^{(a)}(t),\quad
          k^{(b)}(t+dt) \ge k^{(b)}(t).
        \]
  \item \textbf{Интерференция.}\quad
        Межветвевые потоки \(J^{(a\leftrightarrow b)}\) создают
        структурное напряжение, которое подталкивает одну или обе ветви
        ближе к порогам гибели без немедленного схлопывания.
  \item \textbf{Подавление.}\quad
        Одна ветвь переводит другую ниже её порогов жизнеспособности,
        в результате чего подавляемая ветвь схлопывается, тогда как
        подавитель остаётся жизнеспособным или даже усиливается.
  \item \textbf{Паразитизм.}\quad
        Одна ветвь выкачивает поддерживающие потоки у другой, повышая свою
        континуумичность и снижая циклы другой ветви.
  \item \textbf{Слияние.}\quad
        Через сильную связь и общие граничные структуры две ветви могут
        породить новый континуум на общей расширенной ветви. Формально это
        соответствует составному оператору рождения, питающемуся из
        \(B^{(a)}\) и \(B^{(b)}\).
\end{itemize}

Общие граничные структуры особенно важны. Если
\(\partial\Omega^{(a)}\cap\partial\Omega^{(b)}\neq\emptyset\), тогда:
\begin{itemize}
  \item межветвевые потоки, ограниченные пересечением, могут опосредовать
        связь;
  \item состояния на общей границе могут синхронизировать или
        дестабилизировать одновременно обе ветви;
  \item новые континуумы могут зарождаться на общей границе, инициируя новые
        ветви или слияния ветвей.
\end{itemize}

\subsection{Схлопывание ветвей}

Ветви могут завершаться несколькими способами, расширяя понятие
схлопывания от отдельных континуумов
(раздел~the lifecycle and rebirth chapter) до целых траекторий развития.

\paragraph{Определение 13.5 (Схлопывание ветви).}
Ветвь \(B = (K_{x_0}, K_{x_1}, \dots)\) \emph{схлопывается} на уровне
\(x_n\), если:
\begin{enumerate}
  \item континуум \(K_{x_n}\) схлопывается в смысле
        определения~12.1 (пустое допустимое множество, \(k\to 0\),
        нарушение порога);
  \item не существует континуума \(K_{x_{n+1}}\) с \(x_{n+1}>x_n\), на
        котором мог бы быть определён оператор рождения
        \(\Psi_{x_n\to x_{n+1}}\) для какого-либо остатка \(K_{x_n}\)
        (раздел~the lifecycle and rebirth chapter);
  \item ветвь нельзя продолжить дальше при сохранении аксиом размерностной
        монотонности и совместимости вложения.
\end{enumerate}

Структурно различаются три сценария:
\begin{itemize}
  \item \textbf{Терминальное схлопывание.}\quad
        Последний континуум ветви схлопывается и никакого возрождения на более
        высоком уровне не существует. Ветвь конечна и завершается на \(x_n\).
  \item \textbf{Схлопывание через слияние.}\quad
        Ветвь поглощается более устойчивой ветвью через слияние:
        континуумы на \(B\) теряют идентичность как независимые вершины
        графа ветвей и становятся внутренней структурой континуума другой
        ветви. С точки зрения \(B\) это схлопывание через утрату
        идентичности ветви.
  \item \textbf{Асимптотическое истончение.}\quad
        Ветвь не завершается одним резким схлопыванием, а переходит в режим,
        где более высокие уровни становятся всё более предельными
        (\(k_{x_i}\to 0\) при \(i\to\infty\)) и не способны поддерживать
        дальнейшие размерностные рождения. Ветвь бесконечна по индексу,
        но конечна по эффективной сложности.
\end{itemize}

\paragraph{Глобальные ограничения.}
Аксиомы OC накладывают ограничения, предотвращающие неконтролируемые
каскады бесконечного ветвления на конечных уровнях:
\begin{itemize}
  \item пространство вложения на каждом уровне имеет конечную осевую ёмкость для
        данной области, ограничивая число различных каналов рождения;
  \item размерностные пороги требуют нетривиального структурного
        напряжения; многократное ветвление без достаточной опоры потоков
        ведёт ветви к схлопыванию;
  \item теоремы монотонности сложности ограничивают жизнеспособные ветви теми,
        для которых рост сложности совместим с доступными циклами и порогами.
\end{itemize}

\subsection{Последствия для иерархии уровней K}

Структура ветвей уточняет вертикальную иерархию \(K_0\)–\(K_{10}\), определяя,
какие последовательности уровней и осей действительно реализуемы.

\paragraph{Онтогенетические пути развития.}
Каждая ветвь задает \emph{онтогенетический путь развития}:
\begin{itemize}
  \item в физических областях ветви описывают, каким образом континуумы
        теории поля на \(K_2\) порождают химические, протоцеллюлярные и
        биологические континуумы на \(K_3\)–\(K_5\);
  \item в когнитивных и социальных областях ветви описывают, как
        нейроконтинуумы \(K_5\) поддерживают когнитивные континуумы\,
        \(K_6\), которые, в свою очередь, поддерживают социальные и
        цивилизационные континуумы \(K_7\)–\(K_8\);
  \item в теоретических областях ветви соединяют эмпирические континуумы
        с теоретическими и метатеоретическими на уровнях \(K_9\)–\(K_{10}\).
\end{itemize}
Многообразие наблюдаемых физических, химических, биологических и социальных
систем соответствует множественности ветвей, совместимых с общим субстратом.

\paragraph{Допустимые и запрещённые переходы.}
Топология ветвей кодирует, какие переходы между уровнями допустимы:
\begin{itemize}
  \item рёбра графа ветвей \(\mathcal{G}\) соответствуют допустимым
        размерностным операторам рождения;
  \item отсутствие ребра \(K_x \to K_y\) (при \(y>x\)) означает
        структурную невозможность: ни одно пространство вложения и пороговый
        ландшафт не может поддержать такой прямой переход без прохождения
        обязательных промежуточных уровней;
  \item запрещенные переходы обычно нарушают либо размерностную монотонность
        (попытка упрощения осей), либо пороги вложения
        (попытка реализовать оси, неподдерживаемые \(M_x\)).
\end{itemize}

\paragraph{Вертикальная согласованность.}
Поскольку ветви должны редуцироваться к нижним уровням при замораживании
дополнительных осей (раздел~the operator semantics chapter), структура ветвей
обеспечивает вертикальную согласованность:
\begin{itemize}
  \item каждый континуум более высокого уровня \(K_y\) редуцируется к
        континууму нижнего уровня \(K_x\) на ветви при постоянных
        дополнительных осях;
  \item наоборот, нижние уровни выступают необходимыми подструктурами для
        более высоких уровней;
  \item ветви фиксируют эту согласованность как общие префиксы между
        доменами (например, все биологические ветви должны содержать
        физические и химические уровни).
\end{itemize}

\paragraph{Глобальная картина.}
В совокупности вертикальная иерархия \(K_0\)–\(K_{10}\), пороговый
ландшафт и топология ветвей формируют глобальную структурную картину OC.
Операционный аппарат задают объекты \(F\), \(G\), \(H\), \(Q\), \(R\),
\(S\), \(\Psi\), и \(E_{\mathrm{int}}\):
\begin{itemize}
  \item континуумы — локальные объекты с внутренней динамикой;
  \item ветви — глобальные пути развития, проходящие через уровни;
  \item граф ветвей кодирует, какие пути структурно допустимы;
  \item схлопывание и возрождение обрезают и перестраивают ветви со временем.
\end{itemize}

Core~1.3.3 фиксирует эту глобальную картину на структурном уровне. Последующие
расширения и предметно-ориентированные исследования могут квантифицировать
вероятности ветвления, разбирать явные примеры ветвления в конкретных
системах и разрабатывать эмпирические критерии для выделения структуры
ветвей в данных реального мира.

% END OF FILE
